\section*{Hoofdstuk \ref{ch:g8:plastics} --- Kunststoffen en synthetische materialen}

\begin{solution}{exo:g8:plastics:1}
Natuurlijk: katoen, wol. Kunstmatig: papier, viscose. Synthetisch:
nylon, polystyreen.
\end{solution}

\begin{solution}{exo:g8:plastics:2}
Van aardolie (of aardgas).
\end{solution}

\begin{solution}{exo:g8:plastics:3}
PET (polyethyleentereftalaat): water- en frisdrankflessen.
\end{solution}

\begin{solution}{exo:g8:plastics:4}
Een \omterm{def:g8:plastics:thermoplastic}{thermoplast} wordt zacht telkens als je hem verwarmt en kan opnieuw
gevormd worden; een \omterm{def:g8:plastics:thermoplastic}{thermoharder} wordt de eerste keer voorgoed hard en
verkoolt in plaats van zacht te worden.
\end{solution}

\begin{solution}{exo:g8:plastics:5}
Polypropeen (PP), een \omterm{def:g8:plastics:thermoplastic}{thermoplast}: je kunt het smelten en opnieuw
vormen.
\end{solution}

\begin{solution}{exo:g8:plastics:6}
De driehoek geeft alleen de \omterm{def:g8:plastics:plastic}{kunststof} aan, om het sorteren te
vergemakkelijken. Of het voorwerp gerecycled wordt, hangt ervan af of
het ingezameld wordt en of er een bedrijf is dat die \omterm{def:g8:plastics:plastic}{kunststof} recyclet.
\end{solution}

\begin{solution}{exo:g8:plastics:7}
Een pan wordt heet: een \omterm{def:g8:plastics:thermoplastic}{thermoplast} zou zacht worden, een \omterm{def:g8:plastics:thermoplastic}{thermoharder}
niet.
\end{solution}

\begin{solution}{exo:g8:plastics:8}
$\frac{353}{460} \approx 0.77$: ongeveer \qty{77}{\percent}.
\end{solution}

\begin{solution}{exo:g8:plastics:9}
PVC bevat chloor: als het verbrandt, komt waterstofchloride vrij, een
bijtend gas, naast roet en koolstofmono-oxide van de \omterm{def:g8:combustion-and-fuels:complete}{onvolledige
verbranding} in een tuinvuur.
\end{solution}

\begin{solution}{exo:g8:plastics:10}
$300 \times 36 = 10\,800$ flessen; $10\,800 \times 25 = 270\,000$\,g, dat
is \qty{270}{kg}.
\end{solution}

\begin{solution}{exo:g8:plastics:11}
Koolstofdioxide en water. Zijn koolstof komt uit aardolie, die miljoenen
jaren onder de grond opgeslagen was: als je het verbrandt, voeg je
koolstofdioxide toe aan de \omterm{def:g4:air-a-mixture-of-gases:air}{lucht}, net als bij het \omterm{def:g4:what-a-fire-needs:burning}{verbranden} van een
\omterm{def:g8:combustion-and-fuels:fossil}{fossiele brandstof}.
\end{solution}

\begin{solution}{exo:g8:plastics:12}
De fles verrot niet: zon en golven breken haar in steeds kleinere
stukjes. De piepkleine stukjes drijven of zinken in het water, en vissen
slikken ze in met hun voedsel.
\end{solution}

\begin{solution}{pb:g8:plastics:1}
\textbf{1.} Flessen: 1 (PET). Doppen: 2 (HDPE) of 5 (PP).

\textbf{2.} Synthetisch: ze zijn gemaakt van stoffen die scheikundigen
hebben opgebouwd, meestal uit aardolie.

\textbf{3.} Ja: PET is een \omterm{def:g8:plastics:thermoplastic}{thermoplast}.

\textbf{4.} Elke \omterm{def:g8:plastics:plastic}{kunststof} wordt apart gerecycled: als je ze samen
smelt, geeft een \omterm{def:g6:pure-substances-and-mixtures:mixture}{mengsel} van \omterm{def:g8:plastics:plastic}{kunststoffen} een slecht materiaal.

\textbf{5.} $0.80 \times 1000 = \qty{800}{kg}$.

\textbf{6.} Doppen: $0.12 \times 1000 = \qty{120}{kg}$. Etiketten: $0.05
\times 1000 = \qty{50}{kg}$.

\textbf{7.} $0.03 \times 1000 = \qty{30}{kg}$ vuil en vloeistof.

\textbf{8.} $80 + 12 + 5 = \qty{97}{\percent}$.

\textbf{9.} $0.025 \times 800 = \qty{20}{kg}$.

\textbf{10.} De vlokken besparen aardolie en de energie om nieuw PET te
maken, en ze maken van afval een nuttig voorwerp.

\textbf{11.} Koolstofdioxide (en, als de \omterm{def:g4:what-a-fire-needs:burning}{verbranding} slecht is,
koolstofmono-oxide en roet).

\textbf{12.} $800 - 20 = \qty{780}{kg}$ schone PET-vlokken per baal.
\end{solution}
